%
% board_ordering_reread_only_buggy.tex
%
% FIDELITY CONTROL for docs/board_ordering.tex -- candidate fix (b) on its own.
%
% Every push reads the slot at push time rather than pushing a board captured
% earlier -- exactly as the committed design does -- but the writes are not
% serialised.  This variant isolates the gate: it is the committed model with
% the one guard "staged = {}" deleted from PushTake, and nothing else changed.
%
% It answers the question directly: does re-reading the slot close the defect on
% its own, or does the gap between the read and the write still race?  It races.
%
% Type-check:  fuzz -t docs/board_ordering_reread_only_buggy.tex
% Model-check: probcli docs/board_ordering_reread_only_buggy.tex -model_check \
%                -p DEFAULT_SETSIZE 2 -p MAXINT 4 -goal "shown /= evershown"
%              (must report the goal FOUND)
%

\documentclass[a4paper,10pt,fleqn]{article}
\usepackage[margin=1in]{geometry}
\usepackage{fuzz}
\usepackage[colorlinks=true,linkcolor=blue,citecolor=blue,urlcolor=blue]{hyperref}

\begin{document}

\title{Re-reading the Slot Alone: a Fidelity Control}
\author{Why the gap between the read and the write still races}
\date{August 2026}
\maketitle

% ============================================================
\section{Purpose}
% ============================================================

The second candidate fix has the answer read the slot at push time instead of
pushing a \texttt{HeldBoard} captured before the raise. The reasoning is sound
as far as it goes: the board a click shows should be the board the applet holds
when it shows it, not the one it held a round trip ago.

But it treats a wide window as the problem, when the problem is that there is a
window at all. Reading the slot and writing the glass are still two acts with
nothing between them, and two pushers still write in whichever order their
sockets happen to complete in. Narrowing a race does not settle it; it makes
the losing interleaving rarer and therefore harder to reproduce, which is worse
than leaving it obvious.

This variant is \texttt{docs/board\_ordering.tex} with exactly one guard
removed: \texttt{PushTake} no longer requires $staged = \emptyset$. Everything
else --- the state, \texttt{Init}, \texttt{TakePlace}, \texttt{StoreBoard},
\texttt{PushCommit} --- is character for character the committed model. The
difference between a design that holds and a design that does not is that one
guard.

% ============================================================
\section{Types and State}
% ============================================================

\begin{zed}
  [WORKER]
\end{zed}

\begin{zed}
  PLACE == 0 \upto 3
\end{zed}

\begin{schema}{Board}
  began : PLACE \\
  loading : WORKER \pfun PLACE \\
  held : PLACE \\
  shown : PLACE \\
  staged : WORKER \pfun PLACE \\
  everheld : PLACE \\
  evershown : PLACE \\
  offered : PLACE \\
  overtaken : \power PLACE
\where
  held \leq began \\
  shown \leq began \\
  everheld \leq began \\
  evershown \leq began \\
  offered \leq began
\end{schema}

\begin{schema}{Init}
  Board'
\where
  began' = 0 \\
  loading' = \emptyset \\
  held' = 0 \\
  shown' = 0 \\
  staged' = \emptyset \\
  everheld' = 0 \\
  evershown' = 0 \\
  offered' = 0 \\
  overtaken' = \emptyset
\end{schema}

\begin{schema}{TakePlace}
  \Delta Board \\
  w? : WORKER
\where
  w? \notin \dom loading \\
  began < 3 \\
  began' = began + 1 \\
  loading' = loading \cup \{ w? \mapsto began' \} \\
  held' = held \\
  shown' = shown \\
  staged' = staged \\
  everheld' = everheld \\
  evershown' = evershown \\
  offered' = offered \\
  overtaken' = overtaken
\end{schema}

\begin{schema}{StoreBoard}
  \Delta Board \\
  w? : WORKER
\where
  w? \in \dom loading \\
  ( loading~w? > held \land held' = loading~w? \\
  \quad~ \land overtaken' = overtaken ) \lor \\
  \quad~ ( loading~w? \leq held \land held' = held \\
  \quad~ \land overtaken' = overtaken \cup \{ loading~w? \} ) \\
  ( held' > everheld \land everheld' = held' ) \lor \\
  \quad~ ( held' \leq everheld \land everheld' = everheld ) \\
  loading' = \{ w? \} \ndres loading \\
  began' = began \\
  shown' = shown \\
  staged' = staged \\
  evershown' = evershown \\
  offered' = offered
\end{schema}

% ============================================================
\section{The Push, Read Fresh but Unserialised}
% ============================================================

\texttt{PushTake} reads the slot exactly as the committed design does. The one
missing conjunct is $staged = \emptyset$.

\begin{schema}{PushTake}
  \Delta Board \\
  w? : WORKER
\where
  w? \notin \dom staged \\
  staged' = staged \cup \{ w? \mapsto held \} \\
  ( held > offered \land offered' = held ) \lor \\
  \quad~ ( held \leq offered \land offered' = offered ) \\
  began' = began \\
  loading' = loading \\
  held' = held \\
  shown' = shown \\
  everheld' = everheld \\
  evershown' = evershown \\
  overtaken' = overtaken
\end{schema}

\begin{schema}{PushCommit}
  \Delta Board \\
  w? : WORKER
\where
  w? \in \dom staged \\
  shown' = staged~w? \\
  ( staged~w? > evershown \land evershown' = staged~w? ) \lor \\
  \quad~ ( staged~w? \leq evershown \land evershown' = evershown ) \\
  staged' = \{ w? \} \ndres staged \\
  began' = began \\
  loading' = loading \\
  held' = held \\
  everheld' = everheld \\
  offered' = offered \\
  overtaken' = overtaken
\end{schema}

% ============================================================
\section{What it reaches}
% ============================================================

\begin{verbatim}
  probcli docs/board_ordering_reread_only_buggy.tex -model_check \
    -p DEFAULT_SETSIZE 2 -p MAXINT 4 -p TIME_OUT 60000 \
    -goal "shown /= evershown"                          # I2: FOUND
  probcli ... -goal "staged = {} & shown < offered"     # I3: FOUND
  probcli ... -goal "shown > held"                      # I4: not found
  probcli ... -goal "shown : overtaken"                 # I5: not found
  probcli ... -goal "held /= everheld"                  # I1: not found
\end{verbatim}

Two of five found, over 945 reachable states and 2871 transitions, exhaustively
--- and which two is the useful part.

I2 is found, in nine steps: two pushers read the slot, one lands its board, and
the other lands an older one over the top. Re-reading gave the losing pusher a
board that was current when it read it; by the time it wrote, it was not. So (b)
alone does not close the defect --- the gap between the read and the write is
the race, and shortening a race is not settling it.

I4 and I5 are \emph{not} found, and that is what tells us (b) is the right half
of the fix rather than a red herring. Reading the slot at push time is exactly
what stops the glass running ahead of the slot, and exactly what stops a board
the slot refused from ever reaching the screen: a board that comes from the slot
is by definition one the slot holds. The current design violates both; this
variant violates neither. What (b) cannot do on its own is order two writes
against each other. That is what (a) was reaching for and missed, and what the
committed design achieves --- not with a counter, but with the one guard this
variant deletes.

\end{document}
